% !TEX root = ../../tree_notes.tex
% =====================================================================
%  Chapter 3 review -- summary, checklist, quick reference, mixed set,
%  related topics. Uses \chaptersection (unnumbered) and NO numbered
%  boxes (they would continue the 3.4.x numbering).
% =====================================================================

\chaptersection{Chapter 3 Summary}

One question organised this chapter: \emph{when do you write the node down?} The tables
collect the four answers, the facts that make them useful, and the cost of producing each.

\begin{center}
\footnotesize
\renewcommand{\arraystretch}{1.3}
\begin{tabular}{@{}p{0.20\linewidth} p{0.56\linewidth} p{0.18\linewidth}@{}}
\toprule
\textbf{Order} & \textbf{Rule and where the root sits} & \textbf{Defined code} \\
\midrule
Preorder  & node, left, right; root is \textbf{first} -- copy/serialise/paths down & \coderef{c03-dfs3} \\
Inorder   & left, node, right; root \textbf{splits} the list -- BST sorted order & \coderef{c03-dfs3} \\
Postorder & left, right, node; root is \textbf{last} -- delete/evaluate/aggregate & \coderef{c03-dfs3} \\
Level     & by depth, left to right within a level; queue, not stack & \coderef{c03-bfs} \\
\bottomrule
\end{tabular}
\end{center}

\medskip
{\small\textbf{Results of this chapter.}}
\begin{center}
\footnotesize
\renewcommand{\arraystretch}{1.3}
\begin{tabular}{@{}p{0.24\linewidth} p{0.50\linewidth} p{0.20\linewidth}@{}}
\toprule
\textbf{Result} & \textbf{Statement} & \textbf{Where} \\
\midrule
Visit cost        & Every DFS order takes $\bigO{n}$ time, $\bigO{h}$ stack ($2n+1$ calls). & \propref{c03-once-each} \\
Root position     & first in pre, last in post, splitter in inorder. & \propref{c03-root-position} \\
Morris            & $\bigO{n}$ time, $\bigO{1}$ extra space, tree fully restored. & \thmref{c03-morris-cost} \\
BFS orders depths & round $r$ emits exactly depth $r-1$; first visit is nearest. & \propref{c03-bfs-depth} \\
Pre $+$ in unique & distinct values reconstruct exactly one tree. & \thmref{c03-prein-unique} \\
Markers $=$ shape & preorder $+$ nulls round-trips: $2n+1$ tokens. & \coderef{c03-ser} \\
Euler intervals   & $v\in\subtree(u) \iff \tin(u)\le\tin(v)\le\tout(v)\le\tout(u)$. & \propref{c03-euler-interval} \\
\bottomrule
\end{tabular}
\end{center}

\medskip
{\small\textbf{Algorithms from this chapter.}}
\begin{center}
\renewcommand{\arraystretch}{1.2}
\begin{tabular}{@{}p{0.36\linewidth} l l l@{}}
\toprule
\textbf{Algorithm} & \textbf{Time} & \textbf{Memory} & \textbf{Code} \\
\midrule
Recursive pre/in/post      & $\bigO{n}$ & $\bigO{h}$  & \coderef{c03-dfs3} \\
Iterative pre / in / post  & $\bigO{n}$ & $\bigO{h}$  & \coderef{c03-pre-iter}, \coderef{c03-in-iter}, \coderef{c03-post-iter} \\
Morris in (or pre) order   & $\bigO{n}$ & $\bigO{1}$  & \coderef{c03-morris} \\
Level order with splits    & $\bigO{n}$ & $\bigO{w}$  & \coderef{c03-bfs} \\
Rebuild from pre $+$ in    & $\bigO{n}$ & $\bigO{n}$  & \coderef{c03-build-tmpl} \\
Serialise / deserialise    & $\bigO{n}$ & $\bigO{n}$  & \coderef{c03-ser}, \coderef{c03-deser} \\
\bottomrule
\end{tabular}
\end{center}
Here $h$ is the tree height and $w$ its widest level -- the two memory profiles are
independent: skewed trees punish DFS, bushy trees punish BFS.

\begin{figure}[H]
\centering
\begin{tikzpicture}[diagram, node distance=7mm,
    q/.style={boxnode, fill=primaryblue!12, draw=primaryblue, text width=30mm},
    a/.style={boxnode, fill=successgreen!12, draw=successgreen, text width=32mm}]
  \node[q] (ask) {\textbf{What shape does the answer need?}};
  \node[a, below left=10mm and 0mm of ask] (lev) {level by level / nearest nodes \\ \footnotesize BFS (\secref{c03-bfs})};
  \node[a, below=10mm of ask] (ord) {one DFS order \\ \footnotesize children first? postorder;\\ node first? preorder; BST? inorder (\secref{c03-dfs})};
  \node[a, below right=10mm and 0mm of ask] (rec) {rebuild / store a tree \\ \footnotesize pre$+$in (\thmref{c03-prein-unique}) or pre$+$markers (\coderef{c03-ser})};
  \draw[flow] (ask) -- (lev); \draw[flow] (ask) -- (ord); \draw[flow] (ask) -- (rec);
  \node[boxnode, fill=accentorange!12, draw=accentorange, below=8mm of ord, text width=44mm]
    (mem) {\emph{Space pressure?} swap any DFS order for Morris: $\bigO{1}$ (\secref{c03-iterative})};
  \draw[softarrow] (ord) -- (mem);
\end{tikzpicture}
\caption*{Choosing a traversal: breadth vs depth first, which writing moment, or run the
walk backwards into a tree.}
\end{figure}

\chaptersection{Common Mistakes Checklist}

\begin{itemize}[nosep, leftmargin=1.4em]
  \item \textbf{Push order wrong in iterative preorder.} Push \emph{right then left} or the
        walk visits right subtrees first (\coderef{c03-pre-iter}).
  \item \textbf{Freeze too late.} Reading \code{q.size()} \emph{inside} the \code{for} loop
        shuffles level boundaries -- snapshot \code{sz} once per round (\coderef{c03-bfs}).
  \item \textbf{Writing on pop in inorder.} The value must be written \emph{after} the pop
        \emph{and} the left-dive is finished; writing on push gives preorder by accident.
  \item \textbf{Morris forgets to undo.} Remove the thread at the second visit before
        moving right -- leftover threads corrupt the tree for every later algorithm.
  \item \textbf{Serialising without markers.} A bare preorder cannot distinguish the two
        trees of \figref{c03-ambig}; include the \code{\#} nulls (or use two orders).
  \item \textbf{Rebuilding with duplicates or without the hash map.} Values must be
        distinct (\thmref{c03-prein-unique}), and scanning inorder per node gives
        $\bigO{n^2}$ on chains -- index the positions first (\coderef{c03-build-tmpl}).
  \item \textbf{Zigzag parity off by one.} Level $0$ goes left-to-right; test parity with
        the round counter, not the node depth, and check a two-level tree.
  \item \textbf{Confusing level order with ``sorted'' order.} Only \emph{inorder on a BST}
        is sorted (\chref{c05}); level order of a BST is generally not (\figref{c03-runner}
        levels $=[1][2,3][4,5,6]$ happens to be; that is luck).
\end{itemize}

\chaptersection{Quick Reference Card}

\begin{quickref}[Chapter 3 at a glance]
\textbf{The four orders.} pre $=$ NLR, in $=$ LNR, post $=$ LRN, level $=$ by depth.
Root: first / splitter / last / first. Each DFS order is one three-line function -- only
the \code{push\_back} position moves (\coderef{c03-dfs3}).

\medskip
\textbf{Iterative.} Preorder: pop-write, push R then L. Inorder: dive left, pop-write,
go right. Postorder: mirrored preorder into a second stack, pop it reversed.
Morris: borrow null right pointers as temporary threads back to the inorder successor.

\medskip
\textbf{BFS.} \code{sz = q.size()} at the top of each round $=$ one level. Everything
``per level'' (102/103/107/199/515/637) is this loop with one different line.

\medskip
\textbf{Backwards.} pre $+$ in (or post $+$ in) rebuild uniquely; pre $+$ post is
ambiguous without inorder. One order $+$ null markers serialises:
\texttt{1 2 4 \# \# 5 \# \# 3 \# 6 \# \#}.
\end{quickref}

\begin{lstlisting}[caption={Chapter 3 code kit: the two must-memorise walkers side by side}, label={lst:c03-kit2}]
#include <bits/stdc++.h>
using namespace std;
struct TreeNode { int val; TreeNode *left, *right;
                  TreeNode(int x=0): val(x), left(nullptr), right(nullptr){} };

// ITERATIVE INORDER (LeetCode 94 follow-up) -- dive, pop, write, right.
vector<int> inorderIter(TreeNode* root) {
    vector<int> res; stack<TreeNode*> st; TreeNode* cur = root;
    while (cur || !st.empty()) {
        while (cur) { st.push(cur); cur = cur->left; }   // dive all the way left
        cur = st.top(); st.pop();
        res.push_back(cur->val);                          // turnaround moment
        cur = cur->right;
    }
    return res;
}

// LEVEL-ORDER TEMPLATE -- everything "by level" is this loop.
vector<vector<int>> levels(TreeNode* root) {
    vector<vector<int>> res;
    if (!root) return res;
    queue<TreeNode*> q; q.push(root);
    while (!q.empty()) {
        int sz = (int)q.size();                           // freeze the level
        vector<int> lev;
        for (int i = 0; i < sz; i++) {
            TreeNode* v = q.front(); q.pop();
            lev.push_back(v->val);
            if (v->left)  q.push(v->left);
            if (v->right) q.push(v->right);
        }
        res.push_back(lev);
    }
    return res;
}
\end{lstlisting}

\chaptersection{Mixed Problem Set}

The order is not stated -- work out which traversal (and which variant) the task implies.

\begin{mixedset}
  \practice{Validate Preorder Sequence in Binary Search Tree}{LeetCode 331}
    {https://leetcode.com/problems/verify-preorder-sequence-in-binary-search-tree/}{LC Medium}
    {a preorder of a BST obeys a stack of shrinking lower bounds; simulate the walk on the array without building a tree.}
  \practice{Print Binary Tree}{LeetCode 655}
    {https://leetcode.com/problems/print-binary-tree/}{LC Medium}
    {compute the height once, then place nodes on an $h\times 2^h$ grid with the pre-order recursion writing its own coordinates.}
  \practice{Subtree Queries}{CSES 1138}
    {https://cses.fi/problemset/task/1138}{CSES}
    {flatten each subtree to the interval $[\tin,\tout]$ of \propref{c03-euler-interval} and answer sums on the flat array.}
  \practicehh{Vertical Order Traversal of a Binary Tree}{LeetCode 987}
    {https://leetcode.com/problems/vertical-order-traversal-of-a-binary-tree/}{LC Hard}
    {assign each node coordinates (col, row): left child $({}^{-}1,+1)$, right child $(+1,+1)$; group by column, sort inside a column by row then value.}
    {any walk works (BFS keeps rows increasing naturally); collect (row, value) pairs into one vector per column, then sort each vector -- the tie-break rules are the whole problem.}
  \practicehh{Diagonal Traversal of Binary Tree}{GeeksforGeeks}
    {https://www.geeksforgeeks.org/problems/diagonal-traversal-of-binary-tree/1}{GFG Medium}
    {same-diagonal nodes are reached by walking right edges freely and queueing the left child; one queue for the whole traversal.}
    {think of the tree drawn with left edges going down-left: read top-right slant lines one by one; push \code{left} only when the right branch ends.}
\end{mixedset}

\chaptersection{Related Topics}

\begin{itemize}[nosep, leftmargin=1.4em]
  \item \chref{c02} -- the recursion directions these walks embody: preorder $\approx$
        top-down, postorder $\approx$ bottom-up.
  \item \chref{c04} -- arrays and heaps: \emph{storing} the tree so that traversal becomes
        index arithmetic.
  \item \chref{c05} -- binary search trees: inorder sortedness, the reason inorder exists.
  \item \chref{c07} -- n-ary trees: pre/post carry over (589/590/429), inorder does not.
  \item \chref{c08} -- graphs: BFS/DFS with a \code{visited} set; trees are the acyclic
        special case where the parent check suffices.
  \item \chref{c11}, \chref{c16} -- LCA and subtree queries: the $\tin/\tout$ tour teased
        in \secref{c03-rebuild}, now load-bearing.
\end{itemize}
